\subsection{Design composition analysis}
The 24 scale-up annealed designs (Table~\ref{tab:alldesigns}) are all length
25 --- the designer's fixed length --- and their composition is informative
about what the ensemble rewards: cysteine 10.8\%, aromatics F/W 16.0\%,
aliphatics I/V/L 16.8\%, but only 8.0\% K/R. The training positives are 6.8\%
cysteine and 11.9\% K/R, so the designer amplifies cysteine and aromatic
content while actually undershooting the corpus's cationic fraction. Two
readings compete. The charitable one: the ensemble has learned that
Cys-rich, aromatic, hydrophobic short sequences are the most reliably
AMP-like pattern under leakage-controlled training, consistent with the
dominance of disulfide-stabilized families in the reviewed corpus. The
skeptical one: the annealer exploits regions where the model is
overconfident, and cysteine-heavy sequences are exactly where a discriminative
model extrapolates poorly. Distinguishing these requires activity data we do
not have; we therefore present the designs as the model's own optimum, not as
predicted therapeutics.

\subsection{Screen case analysis}
The top screen hits (Table~\ref{tab:allscreen}) share a family resemblance:
short (10--57 aa), cysteine-rich, and hydrophobic, several with apparent
signal-peptide-like N-terminal stretches (e.g.\ the identical P0ACW4/P0ACW5
sequence, 57 aa, 7 Cys). That the same pattern drives both the annealed
designs and the screen ranking is expected --- both consume the same learned
score --- and it means the screen is best read as surfacing small Cys-rich
proteins whose annotation has not caught up with them, rather than as a
diverse portfolio of mechanisms. Two entries in the top 25 (Q54HM1, O55758)
are longer and score lower on the CNN than the GNN relative to their ensemble
rank, illustrating where the two architectures disagree; such disagreement
cases are natural priorities for any follow-up experimental panel, since
consensus at the top is partly an artifact of correlated errors.

\subsection{Reproducibility checklist}
Every claim in this paper is backed by a checked-in artifact:
Tables~\ref{tab:pilot}--\ref{tab:screen} regenerate from
\texttt{results/study01\_amp.json}, \texttt{study07\_scaleup.json},
\texttt{study09\_apd\_validation.json}, and \texttt{study08\_discovery.json}
via \texttt{studies/make\_paper\_assets.py}; Figures~\ref{fig:auroc}
and~\ref{fig:roc} are written by the same script and by study01 respectively;
all data carry provenance manifests (Appendix~\ref{app:prov}); and the test
suite (21 tests) passes at the submission commit. No number in the text was
entered by hand without a corresponding JSON record.
